\begin{theorem}[Alexandrov's theorem for semiconvex functions]
Let $E$ be a finite-dimensional real inner-product space with an additive Haar
measure $\mu$, let $\lambda\ge0$, let $\Omega\subseteq E$ be open and let
$\phi:E\to\R$ be $\lambda$-semiconvex on $\Omega$ in the sense of
\noderef{SemiconvexOn} (so in particular $\Omega$ is convex). Then for
$\mu$-almost every $x\in E$ the following holds: if $x\in\Omega$, there are a
vector $q\in E$ and a symmetric continuous bilinear form $A$ on $E$ with
$A(h,h)\ge-\lambda|h|^2$ for all $h\in E$, such that for every $\delta>0$ there
is $\rho>0$ with
\[
  \Bigl|\phi(x+h)-\phi(x)-\langle q,h\rangle-\tfrac12A(h,h)\Bigr|
    \le\delta|h|^2\qquad(|h|<\rho).
\]
\end{theorem}
\begin{proof}
Set
\[
  \psi(z)=\phi(z)+\frac{\lambda|z|^2}{2}\qquad(z\in E),
\]
so that $\psi$ is convex on $\Omega$ by the definition of semiconvexity in
\noderef{SemiconvexOn}, and $\Omega$ is convex.

\emph{Step 1: a countable cover of $\Omega$ by balls well inside it.}
The space $E$ is a finite-dimensional real normed space, hence separable: fixing
a basis $e_1,\dots,e_n$ of $E$, the countable set of combinations
$\sum_i c_ie_i$ with all $c_i$ rational is dense in $E$. Let
$(x_i)_{i\in\mathbb N}$ enumerate such a countable dense subset of $E$, and let
\[
  \mathcal F=\bigl\{\,(i,s)\ :\ i\in\mathbb N,\ s\in\mathbb Q,\ s>0,\
     \overline B(x_i,2s)\subseteq\Omega\,\bigr\},
\]
a countable set. We claim that the balls $B(x_i,s)$ with $(i,s)\in\mathcal F$
cover $\Omega$. Let $x\in\Omega$. Since $\Omega$ is open there is $r>0$ with
$\overline B(x,5r)\subseteq\Omega$. By density there is $i$ with
$|x_i-x|<r$, and there is a rational $s$ with $r<s<2r$. Then
$|x-x_i|<r<s$, so $x\in B(x_i,s)$; and for $z\in\overline B(x_i,2s)$ we have
$|z-x|\le|z-x_i|+|x_i-x|\le2s+r<4r+r=5r$, so
$\overline B(x_i,2s)\subseteq\overline B(x,5r)\subseteq\Omega$ and
$(i,s)\in\mathcal F$. This proves the claim.

\emph{Step 2: a single null set.}
Let $(i,s)\in\mathcal F$. Applying \noderef{ConvexLocalExtension} to the open set
$\Omega$, the function $\psi$ (convex on $\Omega$), the centre $x_i$ and the
radius $s>0$ --- whose hypothesis $\overline B(x_i,2s)\subseteq\Omega$ holds by
the definition of $\mathcal F$ --- we obtain $\psi_{i,s}:E\to\R$, convex on all
of $E$, with
\begin{equation}\label{eq:ssa-agree}
  \psi_{i,s}=\psi\quad\text{on }B(x_i,s).
\end{equation}
By \noderef{ConvexSecondOrderAe} applied to $\psi_{i,s}$ there is a $\mu$-null
set $Z_{i,s}\subseteq E$ such that every $x\notin Z_{i,s}$ admits a vector
$q\in E$ and a symmetric continuous bilinear form $A$ with $A(h,h)\ge0$ for all
$h$ and
\begin{equation}\label{eq:ssa-exp}
  \Bigl|\psi_{i,s}(x+h)-\psi_{i,s}(x)-\langle q,h\rangle-\tfrac12A(h,h)\Bigr|
   \le\delta|h|^2\qquad(|h|<\rho)
\end{equation}
for every $\delta>0$ and a suitable $\rho=\rho(\delta)>0$. Put
$Z=\bigcup_{(i,s)\in\mathcal F}Z_{i,s}$, a countable union of null sets, hence
$\mu(Z)=0$.

\emph{Step 3: the expansion at every $x\in\Omega\setminus Z$.}
Let $x\in\Omega\setminus Z$. By Step 1 there is $(i,s)\in\mathcal F$ with
$x\in B(x_i,s)$; since $B(x_i,s)$ is open there is $\rho_0>0$ with
$B(x,\rho_0)\subseteq B(x_i,s)$. Since $x\notin Z_{i,s}$, Step 2 provides $q$ and
$A$ with $A(h,h)\ge0$ and \eqref{eq:ssa-exp}.

Let $\delta>0$ and let $\rho=\min(\rho(\delta),\rho_0)>0$. For $|h|<\rho$ both
$x$ and $x+h$ lie in $B(x_i,s)$, so \eqref{eq:ssa-agree} turns
\eqref{eq:ssa-exp} into
\begin{equation}\label{eq:ssa-psi}
  \bigl|e(h)\bigr|\le\delta|h|^2,\qquad
  e(h):=\psi(x+h)-\psi(x)-\langle q,h\rangle-\tfrac12A(h,h)
  \qquad(|h|<\rho).
\end{equation}

Now translate this into a statement about $\phi=\psi-\lambda|\cdot|^2/2$.
Expanding the square,
\[
  \frac{\lambda|x+h|^2}{2}
   =\frac{\lambda|x|^2}{2}+\lambda\langle x,h\rangle+\frac{\lambda|h|^2}{2},
\]
so
\[
  \phi(x+h)-\phi(x)
   =\bigl(\psi(x+h)-\psi(x)\bigr)-\lambda\langle x,h\rangle
     -\frac{\lambda|h|^2}{2}.
\]
Substituting $\psi(x+h)-\psi(x)=\langle q,h\rangle+\tfrac12A(h,h)+e(h)$ from
\eqref{eq:ssa-psi},
\[
  \phi(x+h)-\phi(x)
   =\langle q-\lambda x,h\rangle
    +\tfrac12\bigl(A(h,h)-\lambda|h|^2\bigr)+e(h).
\]
Define
\[
  q'=q-\lambda x,\qquad A'(a,b)=A(a,b)-\lambda\langle a,b\rangle .
\]
Then $A'$ is again a symmetric continuous bilinear form on $E$ (a difference of
two such), $A'(h,h)=A(h,h)-\lambda|h|^2\ge-\lambda|h|^2$ because $A(h,h)\ge0$,
and the display above reads
\[
  \phi(x+h)-\phi(x)-\langle q',h\rangle-\tfrac12A'(h,h)=e(h),
\]
so that $|\phi(x+h)-\phi(x)-\langle q',h\rangle-\tfrac12A'(h,h)|\le\delta|h|^2$
for $|h|<\rho$ by \eqref{eq:ssa-psi}. Since $\delta>0$ was arbitrary and the
pair $(q',A')$ does not depend on $\delta$, the point $x$ admits the asserted
expansion.

As $\mu(Z)=0$, the conclusion holds for $\mu$-almost every $x$: for
$x\notin Z$, either $x\notin\Omega$ and there is nothing to prove, or
$x\in\Omega\setminus Z$ and the expansion has just been established.
\end{proof}
